\section{Compute、Bitter Lesson 与归纳偏置的生命周期}

开场已经给出“识别主导力量”的方法，现在把它用于 AI 研究史。Hyung Won 的核心判断是，计算成本长期下降和 associated scaling 构成重要方向性力量，因此结构设计不能只看今天是否更准，还要看它是否会限制明天继续扩大规模。这里的重点不是宣称单因果，而是把架构选择放到 compute regime 中重新比较。

\subsection{为什么把 compute per dollar 当作主导力量}

本节先读证据。图的纵轴不是晶体管数量，而是固定资金可以买到的计算能力；横轴覆盖超过一个世纪。Hyung Won 由这条长期趋势得出研究策略：不要和强大、持续的外生趋势竞争，应设计能利用它的方法。同时必须把硬件、算法、精度和软件效率都纳入“有效计算”，否则会把某个器件指标错当成整条研究轨迹。

\lecturefigure{hyung-slide-007.jpg}{固定成本可获得的计算能力长期增长，约每五年提高十倍}{Hyung Won Chung official deck, p.~7}

若每五年计算能力约增长十倍，可写成
\[
C(t)=C_0\,10^{t/5},
\]
其中 $C(t)$ 是相对起点经过 $t$ 年后固定预算可购买的计算量，$C_0$ 是起点计算量。这个表达只是在课堂图上的近似，不保证未来继续保持同一指数；能源、制造、互连、资本与算法都可能改变有效增长率。

\begin{knowledgebox}{读图：compute per dollar 不是单一硬件指标}
课堂图聚合了硬件、算法、数值精度、软件栈与系统协同的长期效果。真正相关的问题是“固定成本能完成多少有效训练计算”，而不是只追踪晶体管数量。图展示的是历史趋势；它没有证明能源、制造、互连和资本约束在未来不会改变斜率。
\end{knowledgebox}

\teachervoice{00:08:00--00:09:03 的措辞是“不要与这条趋势竞争，尽量利用它”。这是一条研究选择启发式，不是对无限指数增长的保证；讲义因此把历史观察、研究策略与未来预测分开。}

\subsection{Bitter Lesson：少替机器规定我们自以为懂的思考方式}

上一节说明 scale 是强趋势，本节解释为什么过多人工结构会冲突。研究者常把自己对人类思维的解释编码成模型；这些结构在低算力时可能像捷径，却限制模型探索其他表示。Bitter Lesson 的课堂版本是：长期进展反复来自更一般、假设更弱的方法，再加更多数据与计算。

\lecturefigure{hyung-slide-008.jpg}{“教机器按我们以为自己在思考的方式思考”可能施加错误结构}{Hyung Won Chung official deck, p.~8}

\lecturefigure{hyung-slide-009.jpg}{Bitter Lesson：更一般的方法、更弱的建模假设与更多规模}{Hyung Won Chung official deck, p.~9}

\begin{importantbox}{什么是 inductive bias}
Inductive bias（归纳偏置）是模型在有限数据下偏好某类函数、表示或交互方式的结构性假设。卷积的局部性、encoder/decoder 参数分离、双向 attention、特定不变性都属于偏置。偏置能提高样本效率，但也排除了部分解；它的价值取决于当前数据、算力、任务和部署约束。
\end{importantbox}

\teachervoice{Hyung Won 在 00:09:17--00:11:19 说，研究者往往试图建模“我们以为自己如何思考”，但并不真正理解低层认知机制。课堂观点故意强烈；讲义把它保留为一种设计警告，而不是证明所有先验结构都错误。}

\subsection{结构越少越 scalable，但今天不一定更好}

若把结构多寡放到横向比较，结构较多的方法通常在低计算区间起点更高，却可能较早饱和；结构较少的方法早期训练困难，但给模型更大自由度，随规模增长可能超过前者。这里的 ``scalable'' 指增长曲线能继续改善，不只是能在更多 GPU 上运行。

\lecturefigure{hyung-slide-010.jpg}{更多结构提高低算力性能，却可能降低长期可扩展性}{Hyung Won Chung official deck, p.~10}

\begin{knowledgebox}{读图：比较交叉点，而不是给结构贴好坏标签}
蓝色方法在当前预算左侧可能明显更好；红色方法要到更大 compute 才超越。研究决策取决于部署点在哪里、能否承受探索成本，以及未来是否真的能到达交叉点。图支持的是 regime-dependent choice，而不是“永远选择最少结构”。
\end{knowledgebox}

\lecturefigure{hyung-slide-013.jpg}{当前 compute regime 应选择可工作的偏置，同时记住未来移除它}{Hyung Won Chung official deck, p.~13}

\lecturefigure{hyung-slide-014.jpg}{给定计算、数据、算法和架构时存在暂时最优的 inductive bias}{Hyung Won Chung official deck, p.~14}

可以把结构强度记为 $b$，当前资源状态记为 $r$，性能写成
\[
b^{\star}(r)=\arg\max_b\; P(b;r).
\]
其中 $P(b;r)$ 是偏置 $b$ 在资源状态 $r$ 下的性能，$b^{\star}(r)$ 是当前最优结构。随着 $r$ 包含的 compute、data 与 algorithm 改变，$b^{\star}$ 也会移动；所以一个曾经正确的架构决定不应永久冻结。

\teachervoice{00:12:04--00:13:27 的关键限定是：不能从一开始就等待最一般的方法，因为它可能在当前预算完全不工作。正确动作是“现在加入有用偏置，并明确记录将来何时复查和移除”。}

\subsection{研究激励为什么更擅长加结构，不擅长删结构}

结构生命周期还有社会层原因。新增模块、损失项或机制容易形成论文贡献；证明一个复杂模块在更大规模下可以删除，往往看起来不够新颖。因此社区会积累为旧 regime 设计的假设，却缺少系统性减法。本节把研究激励也视为架构演化的一部分，因为哪些实验得到资源、哪些负结果被发表，会直接影响旧偏置被保留多久。

\lecturefigure{hyung-slide-015.jpg}{长期更好的方法在当前往往更差}{Hyung Won Chung official deck, p.~15}

\begin{warningbox}{短期 leaderboard 优势不等于长期研究价值}
如果一个方法只在小模型、小数据或特定 benchmark 上领先，需要继续问：优势是否来自更强先验？规模扩大后是否饱和？实现复杂度是否妨碍更大训练？这些问题不能由当前单点排名回答。
\end{warningbox}

\teachervoice{Hyung Won 在 00:13:27--00:14:23 直言，论文激励鼓励添加结构，却很少奖励删除结构。他还说长期更好的方法几乎必然在今天看起来更差，因为更自由的学习系统在低资源区间更混乱。}

\lecturefigure{hyung-slide-016.jpg}{阶段总结：主导力量是更便宜的计算与 associated scaling}{Hyung Won Chung official deck, p.~16}

这张总结页不是把其他力量全部否定，而是宣布接下来的分析坐标系：把 Transformer 架构中的额外结构逐项解释为某个历史条件下的 inductive bias，再问当前规模是否仍需要它。这样，架构图不再是静态分类，而变成可以沿 scale trajectory 推理的对象。

\subsection{本章小结}

本章建立了 Hyung Won 的判断框架：计算成本长期下降使一般方法具有强大后劲；归纳偏置在低资源下提高效率，却可能限制后续规模化。正确结论既不是“结构越多越好”，也不是“结构越少越好”，而是持续估计当前 regime 的最优偏置，并为未来的减法实验保留路径。

